学过 VAE、GAN、归一化流和扩散模型之后，脑子里通常留下的是四套互不相干的技巧：一套讲下界与重参数化，一套讲两个网络对着打，一套讲耦合层和 Jacobian 行列式，一套讲加噪声再去噪。四套技巧各有各的词汇表，换一条路线几乎要重学一遍。

它们其实是同一件事的四种做法。生成建模想做的只有一句：让模型分布 \(p_\theta\) 靠近数据分布 \(p_{\text{data}}\)，

\[
\min_\theta\ D(p_{\text{data}}\,\Vert\,p_\theta).
\]

这一行写得出来，却没有一个符号可以直接拿去算。\(p_{\text{data}}\) 那一侧你只有样本，密度值一个也没有；\(p_\theta\) 那一侧更糟——足够灵活的模型通常只写得出一个未归一化的打分，归一化常数是一个高维积分。四条路线的分歧就在这两处：选哪个 \(D\)，以及怎样绕开算不出来的那一项。

选 \(D\) 不是审美问题。散度的方向决定了模型在数据覆盖不到的地方怎么办，而生成模型的成败全在那些地方——本章开篇要讲透的就是这一点：\(\KL(p_{\text{data}}\Vert p_\theta)\) 罚``该有的地方没有''，逼模型摊开去覆盖一切，代价是把质量摊到数据之间的空隙上；\(\KL(p_\theta\Vert p_{\text{data}})\) 罚``不该有的地方有'' ，允许模型缩进少数模式里保证像，代价是别的模式整个消失。样本发糊和模式坍缩这两个经典失败，来源就是这一个不对称，不是两类需要分别调参的工程故障。

绕开归一化常数的四种手法同样可以并排看：VAE 用一个下界避开它，流用可逆性让它恒等于 \(1\)，GAN 根本不建模密度、只学一个采样器，扩散把一次困难的匹配拆成很多次容易的匹配。三样东西——可算的似然、可采样的模型、任意灵活的架构——没有一条路线能同时拿全，选路线本质上是在决定放弃哪一样。

\subsection{怎么读这一章}

第一篇是支点，后三篇都在它建立的坐标上说话，跳过它会把后面的每个结论读成孤立的技巧。第二篇和第三篇是两个极端案例：流为了精确似然把架构锁死，GAN 为了架构自由把密度整个放弃，两笔账正好对照着结。第四篇处理评价，它同时是前三篇的检验——如果第一篇那个不对称是对的，那么质量与多样性就必然是一笔交易，而不是可以同时优化的两个目标。

三条不能弄混的界线：似然高不等于样本好，判别器强不等于生成器有梯度，指标变好不等于模型变强。

\subsection{本章篇目}

\begin{itemize}
\tightlist
\item \hyperref[sec:14-Deep-Learning/07-Generative-Models-Unified/01]{``四条路线在最小化哪个散度''}：两个 KL 方向的惩罚位置为什么相反，模糊与模式坍缩如何从同一个不对称里长出来，四条路线各站在哪个位置，以及归一化常数的四种绕法。
\item \hyperref[sec:14-Deep-Learning/07-Generative-Models-Unified/02]{``归一化流用可逆性换取精确似然''}：变量替换公式为什么把密度和体积变化率绑在一起，维数不变为何不是可以放松的技术细节，耦合层怎样把 \(O(d^3)\) 的行列式变成一次求和，以及精确似然买到了什么、没买到什么。
\item \hyperref[sec:14-Deep-Learning/07-Generative-Models-Unified/03]{``GAN 把散度的估计交给判别器''}：最优判别器的完整推导与它对应的 JS 散度，支撑不重叠时梯度为何消失，Wasserstein 距离与 Lipschitz 约束的代价，以及哪些训练稳定性做法只是经验规律。
\item \hyperref[sec:14-Deep-Learning/07-Generative-Models-Unified/04]{``质量与多样性是一笔交易''}：precision 与 recall 怎样对应两个 KL 方向，似然与感知质量为何能双向背离，采样温度是在权衡曲线上滑动而非提升模型，以及 FID 依赖的三样东西。
\end{itemize}

读完这一章，四条路线之间不再需要各记一套说法：给出训练目标，就能读出它在最小化哪个散度、因此会在哪一侧失败、又是靠什么绕开了不可算的那一项。剩下的共同困难只有一个——训练早期两个分布离得太远，所有散度在那个位置都不好用。下一个单元处理的思路是不去找更好的散度，而是把一次远距离的分布匹配，沿噪声尺度拆成很多次近距离的匹配，每一步都退化成一个平方误差回归。
